%% sections/introduction.tex — Introduction
%% Target: ≥600 words, ~4 paragraphs, ≥8 citations

Understanding how a quantum system restores a broken symmetry after a sudden
perturbation is one of the central questions in non-equilibrium quantum physics.
In many situations of experimental interest, the initial state deliberately breaks
a symmetry that the governing Hamiltonian preserves; the subsequent dynamics then
drives the system toward local symmetry restoration.
How quickly this relaxation proceeds, and whether it can depend counter-intuitively
on \emph{how broken} the initial state was, are questions that bear directly on
the design and interpretation of quantum simulation experiments.

The classical Mpemba effect provides a striking precedent: a hotter sample of
water can freeze faster than a cooler one under identical conditions~\cite{Mpemba1969}.
Although its microscopic mechanism in the classical case remains debated~\cite{Burridge2016},
the effect has inspired a growing body of work on \emph{quantum Mpemba effects}
(QME) — situations where a more strongly out-of-equilibrium quantum state reaches
a target state faster than a less perturbed one.
The entanglement asymmetry (EA), introduced by Ares and collaborators~\cite{Ares2023},
provides a natural quantitative tool for this problem.
The EA of a subsystem $A$ measures how much the reduced density matrix $\rho_A$
departs from the symmetric sector; it vanishes when the subsystem has locally
restored the symmetry.
Using this measure, Ares \textit{et al.}~\cite{Ares2023} demonstrated the quantum
Mpemba effect for U(1) charge symmetry in free-fermion chains initialized in
tilted ferromagnetic states, finding that the more tilted state can restore charge
balance in its subsystem faster than a less tilted one.
This discovery catalysed rapid theoretical activity: quantum Mpemba effects have
since been identified in the anisotropic XY free-fermion chain~\cite{Murciano2024},
open quantum systems~\cite{Carollo2021},
and general integrable systems~\cite{Rylands2024}.

A natural discrete counterpart is the \emph{Z2 quantum Mpemba effect}: whether
a more strongly Z2-broken initial state can restore the Z2 symmetry in a subsystem
faster than a less broken companion.
Two considerations sharpen the importance of this question.
First, in random unitary circuits driven by chaotic dynamics, Z2 QME is expected to be
\emph{absent}: rapid scrambling makes the symmetrisation rate independent of the
initial state, and no Mpemba crossing occurs.
Second, in integrable systems — where long-lived quasiparticle excitations carry
memory of the initial conditions — the Z2 case has remained entirely unexplored.
Given the contrast between U(1) QME (present in integrable free-fermion
chains~\cite{Ares2023}) and Z2 QME (absent in chaotic circuits),
it is unclear whether integrability \emph{enables} Z2 QME through conservation-law-enhanced
memory or whether the discrete symmetry imposes an additional constraint.

A second open question concerns the relationship between QME and another prominent
non-equilibrium phenomenon: dynamical quantum phase transitions (DQPTs).
DQPTs were introduced by Heyl, Polkovnikov, and Kehrein~\cite{Heyl2013} as non-analytic
singularities of the Loschmidt return rate following a quantum quench.
In the transverse-field Ising model (TFIM), DQPTs occur as periodic cusps whenever
the quench crosses the equilibrium phase boundary~\cite{Heyl2013,Heyl2018}.
Because DQPTs represent the most prominent dynamical structure of post-quench evolution
in this model, they suggest themselves as natural candidates for organizing
the timing of Mpemba crossings: perhaps the more broken state restores symmetry
precisely at the moment a DQPT reorganises the quasiparticle vacuum.

In this paper we resolve both questions for the one-dimensional TFIM and its
anisotropic XY extension.
We initialize the system in a product spin state parameterized by a single angle
$\theta$ that continuously controls the degree of Z2 symmetry breaking, and
study the time evolution of the Z2 EA following a sudden quench.
Our main findings are:
\begin{itemize}
  \item \textbf{Z2 QME is robust across the phase diagram.}
    The Z2 quantum Mpemba effect occurs for all 23 post-quench field values tested
    (spanning the FM and PM phases, $\gf \in [0.1, 3.0]$, within the parameter
    range $N \leq 18$, $\NA = 4$), indicating that integrability supports,
    rather than suppresses, Z2 QME.
    This contrasts with the expected absence of Z2 QME in random unitary circuits.

  \item \textbf{QME does not require DQPTs.}
    QME crossings are observed even in the FM phase ($\gf < \gc$), where the
    Loschmidt return rate has no cusps and no DQPT occurs.
    The crossing time $\tM$ evolves continuously through the quantum phase transition
    at $\gc$ without any anomaly.

  \item \textbf{The initial EA imbalance governs $\tM$, not DQPT timing.}
    Across 671 parameter pairs, $\tM$ correlates strongly with the initial EA
    imbalance $|\Delta\Delta\mathcal{S}|$ (Spearman $\rho = +0.90$,
    $p \approx 3 \times 10^{-240}$).
    The distribution of $\tM$ within each DQPT period is significantly non-uniform
    (KS $p \approx 6 \times 10^{-25}$), but the non-uniformity is not centred
    on the DQPT location (binary test $p = 0.27$): DQPTs modulate but do not
    trigger crossings.

  \item \textbf{The effect is robust.}
    Z2 QME persists in the thermodynamic limit (finite-size scaling over $N = 8$--$18$),
    survives variation of the subsystem size ($\NA = 2$--$6$), and is present
    across the anisotropic XY chain family ($\gamma = 0$ to $\gamma = 1$) with
    crossing-time variations of at most $8\%$.
\end{itemize}

An exact closed-form expression for the initial EA,
$\DS_A(0;\theta) = \hbin\!\bigl((1 + \cos^{\NA}\!\theta)/2\bigr)$,
anchors these results and provides an unambiguous, machine-precision-verified
characterisation of the initial conditions.
This formula shows that $\DS_A(0)$ is bounded by $\ln 2$ for any Z2 symmetry
and is controlled entirely by the initial polarisation angle $\theta$ — a parameter
that is directly accessible in cold-atom and quantum simulation platforms~\cite{Bloch2008}.

The remainder of the paper is structured as follows.
Section~\ref{sec:TFIM} defines the model, initial state, and observables.
Section~\ref{sec:DS0}--\ref{sec:xy} presents the numerical results.
Section~\ref{sec:discussion} discusses the physical mechanism, the contrast with
chaotic dynamics, and open directions including analytical understanding and
experimental realisation.
